\clearpage
\section{The complete cubic-background formula}
\label{sec:cubic}

We now determine the entire projection with highest retained derivative
index three. The words ``cubic background'' refer to the fact that a
formal shift of $u_0$ is a cubic polynomial when $u_r=0$ for $r>3$.
They do not restrict the power of $u_3$ in the answer.

\begin{theorem}[Complete cubic normal form]\label{thm:main}
With the notation of \eqref{eq:basic-variables},
\begin{equation}\label{eq:main-generating}
 \sum_{m\ge0}\sigma_3N_m\frac{t^m}{m!}=e^{u_0t}F(y,z,w),
\end{equation}
where
\begin{equation}\label{eq:main}
\boxed{\begin{aligned}
 F(y,z,w)={}&f_{\mathrm{aff}}(y)
 +\frac{z^2}{192}\avg{P_2e^{yv}\Erem(zh)}\\
 &-\frac{w^2}{9216}\avg{P_3e^{yv+zh}\Erem(wk)}.
\end{aligned}}
\end{equation}
This function is entire jointly in $y,z,w$. Writing
$F=\sum_{n\ge0}w^nF_n(y,z)$, its coefficients are $F_0$ from
Theorem~\ref{thm:quadratic}, $F_1=0$, and
\begin{equation}\label{eq:Fn-density}
 F_n(y,z)=-\frac1{9216n!}
 \avg{P_3(v)k(v)^{n-2}e^{yv+zh(v)}},\qquad n\ge2.
\end{equation}
\end{theorem}

\subsection{The closed Euler equation}
Set
\[
 A=2y+\frac z2\partial_y,\qquad B=\partial_z,\qquad
 C=\frac12\partial_y,\qquad L=A+wB.
\]
On a function $e^{u_0t}G(y,z,w)$ the derivation $\delta_3$ acts as
$e^{u_0t}t^{-1}LG$. Lemma~\ref{lem:finitejet} therefore gives
\begin{equation}\label{eq:cubic-PDE}
 F-\frac12L F_y+L^2F_z-L^3F_w
 =\avg{e^{yv+zh(v)+wk(v)}}.
\end{equation}
In detail, after $s=t(1+v)/2$ in \eqref{eq:Euler-A}, the exponent of
the shifted Bell exponential is
\[
 u_0t+\frac{u_1t^2v}{2}
 +\frac{u_2t^3(1+3v^2)}{24}
 +\frac{u_3t^4(v+v^3)}{48}.
\]
This establishes the right-hand side of \eqref{eq:cubic-PDE} with all
normalizations fixed.

The commutators $[B,A]=C$ and $[C,A]=1$ give
\begin{align*}
 L^2&=A^2+w(2AB+C)+w^2B^2,\\
 L^3&=A^3+w(3A^2B+3AC+1)
          +w^2(3AB^2+3CB)+w^3B^3.
\end{align*}
Put
\[
 R_n(y,z)=\frac1{n!}\avg{k(v)^ne^{yv+zh(v)}}.
\]
Comparison of $w^n$ in \eqref{eq:cubic-PDE} gives, for $n\ge1$,
\begin{equation}\label{eq:cubic-rec}
\begin{split}
 (n+1)A^3F_{n+1}={}&(1-n)F_n-\frac{1+3n}{2}A\partial_yF_n
 +(1-3n)A^2\partial_zF_n\\
 &+(5-3n)A\partial_z^2F_{n-1}
 +\frac{3(1-n)}2\partial_y\partial_zF_{n-1}\\
 &+(3-n)\partial_z^3F_{n-2}-R_n.
\end{split}
\end{equation}
Negative subscripts mean zero. The $w^0$ equation is exactly the
quadratic equation because the normal basis excludes a single $u_3$.

\begin{lemma}[Uniqueness of the cubic hierarchy]\label{lem:injective}
The operator $A^3$ is injective on $\C[[y,z]]$. Consequently
\eqref{eq:cubic-rec}, together with the specified $F_0$ and $F_1=0$,
determines every $F_n$ uniquely.
\end{lemma}
\begin{proof}
If $G$ is nonzero, let $j$ be its least nonzero $z$-degree and write that
coefficient as $g(y)$. The coefficient of $z^j$ in $A^3G$ is
$8y^3g(y)$, which is nonzero. Indeed every use of $(z/2)\partial_y$
instead of $2y$ raises $z$-degree. This proves injectivity and then the
successive uniqueness assertion.
\end{proof}

\subsection{Why the beta density solves every recurrence}
Let $E(v)=e^{yv+zh(v)}$. The key transport relation is
\begin{equation}\label{eq:transport}
 A E(v)=2\partial_v E(v),\qquad
 A^j\avg{Q(v)E(v)}=(-2)^j\avg{Q^{(j)}(v)E(v)}.
\end{equation}
The second equality holds when $Q,Q',\ldots,Q^{(j-1)}$ vanish at
both endpoints. For $Q$ divisible by $P_3$, it applies for $j\le3$.
The coefficients of $A$ are independent of $v$, so iteration of the
first equality and ordinary integration by parts prove the second.

\begin{proof}[Proof of Theorem~\ref{thm:main}]
We verify that \eqref{eq:Fn-density} satisfies
\eqref{eq:cubic-rec}. Uniqueness then identifies it with the actual
normal projection, so no unproved assumption about a formal solution is
needed.

For $n=1$, the recurrence reduces to
\[
 2A^3F_2=2A\partial_z^2F_0-R_1.
\]
Using \eqref{eq:F0zz} and \eqref{eq:transport}, this is the polynomial
identity
\begin{equation}\label{eq:cubic-base1}
 \frac{P_3'''}{1152}=-\frac{P_2'}{48}-k.
\end{equation}
For $n=2$, the required identity is
\begin{equation}\label{eq:cubic-base2}
 -8(P_3k)'''=-P_3+7(vP_3)'-20(hP_3)''-96P_2h+9216k^2.
\end{equation}
Both follow by differentiating the displayed polynomials.

For every $n\ge3$, substituting \eqref{eq:Fn-density} in the
recurrence, multiplying by $-9216n!$, and using
\eqref{eq:transport} reduces the assertion to
\begin{equation}\label{eq:cubic-density-identity}
\begin{split}
 -8(P_3k^{n-1})'''={}&(1-n)P_3k^{n-2}
 +(1+3n)(vP_3k^{n-2})'\\
 &+4(1-3n)(hP_3k^{n-2})''
 -2n(5-3n)(h^2P_3k^{n-3})'\\
 &+\frac32n(1-n)vhP_3k^{n-3}
 +n(n-1)(3-n)h^3P_3k^{n-4}\!\\
 &+9216k^n.
\end{split}
\end{equation}
At $n=3$ the term with factor $3-n$ is zero; no negative power is
part of the asserted identity. For $n\ge4$ all terms are polynomials.

Here is a finite algebraic way to verify \eqref{eq:cubic-density-identity}
for an indeterminate $n$, making the all-index step explicit. Define
\[
 \mathcal G_{a,q}(U)=U'+a\frac{q'}qU.
\]
On the open set where $k\ne0$, divide \eqref{eq:cubic-density-identity}
by $k^{n-4}$ and replace its left side by
$-8\mathcal G_{n-4,k}^3(P_3k^3)$. The terms on its right become,
in order,
\begin{align*}
 &(1-n)P_3k^2,
 &&(1+3n)\mathcal G_{n-4,k}(vP_3k^2),\\
 &4(1-3n)\mathcal G_{n-4,k}^2(hP_3k^2),
 &&-2n(5-3n)\mathcal G_{n-4,k}(h^2P_3k),\\
 &\tfrac32n(1-n)vhP_3k,
 &&n(n-1)(3-n)h^3P_3,
 \qquad 9216k^4.
\end{align*}
Substitute $P_3=(1-v^2)^3$, $24h=1+3v^2$, and $48k=v+v^3$.
Expanding the derivatives and collecting powers of $n$ makes the
difference identically zero in $\Q(v)[n]$. Only three derivatives are
needed, independently of $n$. Clearing denominators proves the identity
on the entire polynomial domain. The same calculation at $n=3$, or its
polynomial continuation, proves that case with the stated zero-term
convention. The supplied script records these exact symbolic operations.

All integrations by parts in this verification have zero endpoint terms:
the worst derivative transfer is of order three, and its density has the
factor $P_3$. The exceptional terms involving $F_0$ occur only at
$n=1,2$ and were handled separately. Thus the recurrence is verified at
every index. Lemma~\ref{lem:injective} proves
\eqref{eq:Fn-density}. Finally,
$\sum_{n\ge2}w^nk^{n-2}/n!=w^2\Erem(wk)$ gives
\eqref{eq:main}. The integrands are entire and uniformly bounded on
compact parameter sets and the compact interval, proving joint entirety.
\end{proof}

\begin{remark}[Proof status of the algebraic step]
The four identities \eqref{eq:pure-density-identity},
\eqref{eq:cubic-base1}, \eqref{eq:cubic-base2}, and
\eqref{eq:cubic-density-identity} are explicit finite polynomial or
rational-function identities. Their symbolic verification treats $n$ as
an indeterminate. The separate checks at weights at most 30 are useful
regression tests but are not the reason that the theorem holds for every
weight.
\end{remark}
